\documentclass[12pt]{article}

\usepackage[utf8]{inputenc}
\usepackage{amsmath, amssymb, amsthm}
\usepackage{geometry}
\geometry{a4paper, margin=2.5cm}

\title{Teorema Graficului Închis}
\author{}
\date{}

\theoremstyle{definition}
\newtheorem{definition}{Definiție}

\theoremstyle{plain}
\newtheorem{theorem}{Teoremă}
\newtheorem{proposition}{Propoziție}

\begin{document}

\maketitle

\begin{definition}
Fie $X$ și $Y$ două spații Banach. Pentru un operator liniar $T : X \to Y$, \emph{graficul} lui $T$ este mulțimea


\[
\Gamma(T) = \{ (x, Tx) : x \in X \} \subset X \times Y.
\]


Spunem că $T$ are \emph{grafic închis} dacă $\Gamma(T)$ este închisă în spațiul produs $X \times Y$ (cu norma
$\|(x,y)\| = \|x\| + \|y\|$).
\end{definition}

\begin{theorem}[Teorema Graficului Închis]
Fie $X$ și $Y$ două spații Banach. Dacă $T : X \to Y$ este un operator liniar cu graficul închis, atunci $T$ este continuu.

Echivalent, dacă din


\[
x_n \to x \quad \text{în } X, \qquad Tx_n \to y \quad \text{în } Y
\]


urmează că $y = Tx$, atunci $T$ este continuu.
\end{theorem}

\begin{proposition}
Pentru un operator liniar $T : X \to Y$ între spații Banach, următoarele afirmații sunt echivalente:
\begin{enumerate}
    \item[(i)] $T$ este continuu.
    \item[(ii)] Graficul $\Gamma(T)$ este închis în $X \times Y$.
    \item[(iii)] Există o constantă $C > 0$ astfel încât
    

\[
    \|Tx\| \le C \|x\| \quad \text{pentru toate } x \in X.
    \]


\end{enumerate}
\end{proposition}

\end{document}
